\section{Fase 1: Desarrollo del entorno de autoría}
\label{sec:Fase1Resumen}

La primera fase del proyecto tuvo como objetivo desarrollar una
herramienta gráfica que permitiera crear los cartuchos del sistema de
forma visual, evitando la edición manual de ficheros de configuración
complejos. Esta herramienta constituye el punto de partida de todo el
ecosistema, ya que es la encargada de generar la información utilizada
posteriormente por los motores de simulación y ejecución.

Su desarrollo se organiza en dos bloques claramente diferenciados: la
\textbf{base} —la carga del \emph{setup} que describe la máquina física
y los formatos de cartucho que viajan entre las etapas del sistema— y
el \textbf{diseñador gráfico}, la herramienta de edición visual que es
el corazón de la fase y el auténtico generador de datos del proyecto.

\subsection{Objetivos de la fase}
\label{sec:F1Objetivos}

Los objetivos principales fueron:

\begin{itemize}
	\item Cargar y validar la configuración hardware de una máquina
	Game \& Watch.
	\item Cargar un cartucho definido mediante imágenes y ficheros JSON.
	\item Extraer automáticamente los segmentos gráficos a partir de una
	imagen fuente.
	\item Permitir definir visualmente propiedades y relaciones entre
	segmentos.
	\item Generar un formato de cartucho intermedio reutilizable por el
	simulador.
	\item Facilitar la modificación posterior de los cartuchos mediante
	mecanismos de edición, guardado y recarga.
\end{itemize}

\subsection{La base: el setup de la máquina y los formatos de cartucho}
\label{sec:F1Base}

El punto de partida del entorno es la carga de dos tipos de información
(las primeras iteraciones de la fase): un fichero de configuración
general (\textbf{setup}) que describe las características de la máquina,
y un \textbf{cartucho} que contiene las imágenes y parámetros
específicos del juego. Sobre ellos se establecieron los mecanismos
básicos de validación, detección de errores y presentación de mensajes
al usuario.

\textbf{Decisiones de diseño y arquitectura.}

\begin{itemize}
	\item \textbf{Aislamiento del entorno local:} se prohibió el uso de
	módulos ECMAScript (import/export) en favor de scripts clásicos con un
	ámbito global compartido. Esta decisión neutralizó las restricciones
	del modelo CORS, permitiendo ejecutar el diseñador abriendo
	directamente el archivo \texttt{index.html} mediante el protocolo
	\texttt{file://}, sin servidores de compilación.
	\item \textbf{Abstracción de estilos de máquina:} los botones de
	hardware analizados se tipificaron en una matriz de tres estructuras
	fijas (Estilos 2, 4 y 5). El motor altera el árbol del DOM y traslada
	el lienzo de dibujo al hueco exacto del esqueleto de cada estilo sin
	provocar re-flujos visuales.
	\item \textbf{Tres estados del cartucho:} se definió el ciclo de vida
	del JSON en tres estados de representación: \textbf{GWB} (Base, solo
	geometría), \textbf{GWI} (Intermedio, fuente de verdad de la geometría
	con la lógica acoplada) y \textbf{GWC} (Compilado final binario
	optimizado para el firmware del ESP32).
	\item \textbf{Determinismo indexado:} para que los datos de autoría
	sobrevivan a futuras alteraciones de la imagen, los cambios manuales
	se registran vinculados a los índices primitivos de las componentes
	extraídas, estrictamente estables mientras la geometría base no
	sufra modificaciones por píxeles.
\end{itemize}

\textbf{Problemas detectados y soluciones implementadas.}

\begin{itemize}
	\item \textbf{El problema de la inconsistencia en la botonera:} el
	intento de instanciar la botonera directamente desde la configuración
	de hardware generaba desbordamientos visuales cuando el cartucho
	declaraba un subconjunto de controles inferior al del estilo base de
	la máquina.
	\item \textbf{La solución de diseño:} un sistema de validación
	bidireccional con destrucción selectiva del DOM. El subsistema
	verifica la restricción de compatibilidad (cartucho $\leq$ setup),
	renombra dinámicamente las etiquetas físicas según el identificador de
	autoría y elimina del árbol los nodos sobrantes, garantizando que el
	aparato permanezca libre de referencias colgadas.
	\item \textbf{El problema de la escritura en disco:} el modelo de
	aislamiento (\emph{sandbox}) de los navegadores modernos bloqueaba el
	guardado asíncrono mediante la API de acceso al sistema de archivos
	en entornos no basados en Chromium, impidiendo la descarga autónoma
	del proyecto.
	\item \textbf{La solución de diseño:} se integró la librería de
	compresión JSZip inyectada como script tradicional. El guardado
	empaqueta el manifiesto (una línea compacta por elemento) junto con
	las imágenes de los fragmentos y los lienzos originales, y al
	finalizar crea un nodo temporal de anclaje de descarga en el DOM que
	fuerza de forma universal la descarga del archivo
	\texttt{cartucho\_gwi.zip}.
\end{itemize}

	\subsection{Formato del fichero de setup}

En este fichero se define la máquina que vamos a usar a nivel de hardware.
Se puede disponer de varios ficheros de configuración de tipo \textbf{json}, dependiendo de los elementos existentes y sus conexiones en la máquina de destino. 
\clearpage
\begin{lstlisting}[
	language=json, 
	frame=single, 
	breaklines=true, 
	basicstyle=\ttfamily\small,
	caption={Estructura del archivo setup.json},
	label={lst:setup_base}
	]
	{
		'tipo': 'GW',                                       
		'version': '1.1',                                   
		'nombre': 'ESP32 GW con 5 botones',                 
		'estilo': '5',                                      
		'pantalla': { 'ancho': 320, 'alto': 240 },          
		'botones': [                                        
		{ 'numero': '0', 'puerto': 13, 'tecla': 'Q' },  
		{ 'numero': '1', 'puerto': 14, 'tecla': 'P' },  
		{ 'numero': '2', 'puerto': 25 },                
		{ 'numero': '3', 'puerto': 26 },                
		{ 'numero': '4', 'puerto': 27 },
		{ 'numero': '5', 'puerto': 32 },
		{ 'numero': '6', 'puerto': 33, 'alias': 'A' },
		{ 'numero': '7', 'puerto': 16, 'alias': 'B' },
		{ 'numero': '8', 'puerto': 17, 'alias': 'R' }
		]
	}
\end{lstlisting}

\subsection{El diseñador gráfico: extraer la imagen maestra}
\label{sec:F1Extraccion}

En lugar de obligar al creador del juego a importar decenas de
archivos PNG individuales, el diseñador ingiere una única imagen
maestra con todos los elementos activados a la vez
(\texttt{referencia.png}) y localiza y aísla cada geometría de manera
automática.

\textbf{Decisiones de diseño y arquitectura.}

\begin{itemize}
	\item \textbf{Modelado por componentes conexas:} el algoritmo de
	inspección implementa una búsqueda en anchura (BFS) optimizada
	mediante una pila nativa, evaluando el mapa de píxeles con
	conectividad-8, de modo que los trazos diagonales continuos se
	registran como una misma entidad visual.
	\item \textbf{Deduplicación píxel-exacta:} para mitigar la huella de
	memoria de cara al hardware embebido, el software calcula una firma
	de igualdad lógica en tres dimensiones estrictas: ancho, alto y
	coincidencia binaria de la máscara de píxeles. Las geometrías
	idénticas que se repiten en pantalla se almacenan una única vez en la
	tabla de \emph{sprites}, mientras que sus posiciones individuales se
	segregan en la tabla de instancias.
\end{itemize}

\textbf{Problemas detectados y soluciones implementadas.}

\begin{itemize}
	\item \textbf{El problema del ruido geométrico en el lienzo:} los
	restos accidentales de las herramientas de dibujo digital (píxeles
	sueltos, suavizados involuntarios en los contornos, motas
	imperceptibles) eran detectados por el motor BFS como segmentos
	independientes válidos, generando decenas de archivos fantasma que
	colapsaban la memoria dinámica y corrompían el orden determinista de
	la numeración.
	\item \textbf{La solución de diseño:} un filtro de descarte de ruido
	por umbral de masa geométrica. Cualquier componente conexa con un área
	estrictamente inferior a 4 píxeles reales se intercepta en tiempo de
	ejecución, se purga de las colas de asignación y se notifica con un
	mensaje de advertencia en la consola de depuración para su subsanación
	en el origen gráfico: nada falla en silencio.
\end{itemize}

\subsection{El primer gran replanteamiento: la semántica abandona la imagen}
\label{sec:F1Replanteo}

La propuesta original para dotar de comportamiento lógico a las
geometrías extraídas consistía en codificar parte de la semántica del
juego en la propia imagen origen mediante una convención de colores:
negro para los segmentos normales, verde para el jugador, azul para
las secuencias de movimiento y rojo para las relaciones especiales.
La solución parecía atractiva porque describía gran cantidad de
información con una única imagen; el extractor llegó a realizar un
censo cromático previo que exponía la paleta detectada para que el
usuario la mapeara con los roles del sistema antes de procesar.

\textbf{Problemas detectados y soluciones implementadas.}

\begin{itemize}
	\item \textbf{El problema del colapso cromático por anti-aliasing:}
	las pruebas de estrés con imágenes generadas por editores
	convencionales arrojaron 119 tonos distintos en un lienzo diseñado
	conceptualmente con solo 8 colores puros. Los algoritmos invisibles de
	suavizado periférico de los programas comerciales destruían la
	conectividad exacta del modelo, rompiendo secuencias y uniones de
	manera silenciosa. El comportamiento hacía imposible garantizar una
	interpretación consistente de la imagen en todos los casos.
	\item \textbf{La solución de diseño (primer gran replanteamiento):} se
	determinó que delegar la precisión semántica en un editor de imágenes
	de terceros constituía un defecto de diseño arquitectónico. Se
	rediseñó el flujo por completo: las imágenes quedaron limitadas
	estrictamente a la información geométrica en blanco y negro, y toda
	la lógica de relaciones, secuencias y roles fue trasladada al dominio
	exclusivo del diseñador gráfico mediante operaciones visuales de
	selección. La información gráfica y la información semántica quedan
	desde este hito separadas: las imágenes representan; el editor
	significa.
\end{itemize}

Esta decisión es la que convierte al diseñador gráfico en la pieza
más fundamental del ecosistema: si toda la semántica vive en él, es
el editor —y no un artefacto externo— quien garantiza la coherencia
de los cartuchos.

\subsection{El diseñador gráfico: edición visual de la semántica}
\label{sec:F1Editor}

Sobre el nuevo modelo, las iteraciones restantes de la fase
construyeron la herramienta de edición: paneles de propiedades y
relaciones por segmento, unión de fragmentos, la visualización
simultánea del lienzo de trabajo y de la pantalla real del aparato, y
el mecanismo de deshacer/rehacer.

\textbf{Decisiones de diseño y arquitectura.}

\begin{itemize}
	\item \textbf{Flujo cronológico atómico:} se sustituyeron las matrices
	de estados dispersas por una única lista lineal de operaciones
	atómicas. Cada acción del usuario se inserta con un identificador de
	operación correlativo; si una manipulación provoca múltiples
	mutaciones en cascada, todos los registros comparten el mismo
	identificador de enlace, y deshacer y rehacer revierten o reponen el
	grupo completo sin corromper la integridad del formato intermedio.
	\item \textbf{Visualización transitiva y direccional:} el mapa de
	edición traza las relaciones por color de selección: el componente
	activo en rojo, sus cadenas salientes en azul (calculadas de forma
	transitiva hasta agotar la ruta) y los enlaces entrantes directos en
	magenta.
	\item \textbf{Renderizado fiel al cristal líquido:} el motor gráfico
	de la autoría se independizó como lector puro de la lista de memoria,
	con repintado completo por cambio de estado (fondo estático más
	segmentos activos) y escalado por vecino más cercano, que estira las
	imágenes al vuelo según la escala por eje del simulador sin difuminar
	los contornos del LCD original.
\end{itemize}

\textbf{Problemas detectados y soluciones implementadas.}

\begin{itemize}
	\item \textbf{El problema de la colisión de cajas en el trazado de
		sprites:} las cajas envolventes de los segmentos son rectángulos
	rígidos, pero las formas físicas del cristal líquido son irregulares y
	curvas. Al pintar un segmento con una máscara rectangular sobre una
	zona adyacente, el recuadro opaco borraba los detalles gráficos de
	los elementos que residían inmediatamente debajo.
	\item \textbf{La solución de diseño:} se implementó la mezcla por
	operación binaria XOR canal a canal sobre el vector de bytes del
	lienzo. Los píxeles oscuros se combinan algebraicamente con el fondo
	manteniendo la transparencia absoluta de los píxeles blancos, y el
	borrado exacto de una geometría se ejecuta repitiendo la misma
	función con idénticos argumentos, restituyendo el fondo sin necesidad
	de mapas de bits de respaldo en memoria.
	\item \textbf{El problema de las referencias colgantes por
		desaparición de segmentos:} al unificar múltiples elementos en un
	único sprite compuesto, el componente seleccionado absorbía a los
	secundarios y los enlaces lógicos (\emph{anterior}/\emph{siguiente})
	de otros segmentos quedaban apuntando a identificadores numéricos
	obsoletos que ya no existían en el mapa, colapsando el bucle del
	futuro simulador.
	\item \textbf{La solución de diseño:} se reestructuró la función de
	fusión (\emph{merge}) como una transacción de base de datos: al
	consolidarse una unión, un barrido automático recorre la totalidad de
	la lista y re-apunta cualquier enlace externo que hiciese referencia a
	un miembro absorbido, redirigiéndolo al identificador del nodo raíz
	resultante. Los elementos anulados se transforman en nodos
	informativos fantasma (tipo 99) para preservar estables los índices
	deterministas de la extracción primaria.
\end{itemize}

\subsection{Cierre de la fase}
\label{sec:F1Cierre}

Con la arquitectura del editor congelada, la fase se cierra con la
herramienta de autoría completa: el \emph{setup} describe la máquina,
la imagen maestra aporta la geometría, y el diseñador gráfico —
mediante operaciones visuales, historial atómico y validaciones que
nunca fallan en silencio— produce el GWI, el formato intermedio que
es la fuente de verdad de todo el sistema y el punto de partida de la
fase siguiente: el motor gráfico del simulador.
\clearpage